A detector scoring 99\% on a machine-generated-text benchmark may have learned how machine language differs from human language. It may instead have learned how one corpus was punctuated. The headline number does not separate the two. We decompose the signal into three arms. A surface-only model reads 47 orthographic features and never sees a word. A content-only model reads text stripped of punctuation, casing and non-ASCII characters. A fine-tuned transformer on raw text is the reference. On HC3 the two arms are indistinguishable under a paired test, 3.20\% against 3.26\% error, and the null holds on five of five partitions with the sign changing between them. DeBERTa reaches 0.28\% on the same split, so the parity joins two weak arms rather than a ceiling. On DAIGT V2 content is stronger by a factor of 7.9 in error. Subgroup splits localise that result rather than confirm it. The surface arm reaches 0.01\% error on the Reddit sub-corpus, which is 74.8\% of HC3, and loses to content on both other domains large enough to test. Removing the cue responsible, the rate of spaces before punctuation, costs that arm 10.03 points. BERT emits identical token identifiers for strings differing only by that cue, so it cannot represent it, and it still reaches 0.9916 weighted F1. Surface form is informative on both corpora. We claim only that it reaches parity with content on one, and there for one reason.
